%! The Quadratic Formula & the Discriminant — Unit 2, Lesson 2.5 — Guided Notes
\documentclass[10pt]{article}
\usepackage{algebra2-article}
\usepackage{algebra2-boxes}
\newcommand{\TallMath}[1]{$\displaystyle #1\rule[-1.4em]{0pt}{3.2em}$}
\newcommand{\lgrid}[4]{%
  \draw[step=1, linegray!40, very thin] (#1,#3) grid (#2,#4);
  \draw[->, semithick, charcoal] (#1,0)--(#2,0) node[below right, font=\scriptsize]{$x$};
  \draw[->, semithick, charcoal] (0,#3)--(0,#4) node[left, font=\scriptsize]{$y$};}

\begin{document}

\pageheader{Unit 2, Lesson 2.5}{Guided Notes}

\vspace{0.10in}

\begin{objectivebox}
\textbf{By the end of this lesson, I will be able to\dots}
\begin{itemize}\setlength{\itemsep}{2pt}
  \item solve \emph{any} quadratic with the \emph{quadratic formula}, giving real or complex answers in simplest form;
  \item compute the \emph{discriminant} $b^2-4ac$ and predict the number and type of roots \emph{before} solving;
  \item choose an efficient method, verify a solution, and model a situation with a quadratic.
\end{itemize}
\end{objectivebox}

\vspace{0.06in}

\begin{vocabbox}
\small
Fill in each term as we build it together.
\par\vspace{2pt}
\termblanklong{Quadratic formula}
\termblanklong{Standard form $ax^2+bx+c=0$}
\termblanklong{Discriminant $b^2-4ac$}
\termblanklong{Double (repeated) root}
\termblanklong{Complex-conjugate roots}
\end{vocabbox}

\begin{hookbox}
Four quadratics, four different ``best methods'': $x^2-2x-3=0$ (factor), $x^2-6x+7=0$ (won't factor),
$x^2+2x+5=0$ (complex roots), $2x^2-7x+3=0$. Choosing the right tool each time is a hassle. \textbf{What if
one formula solved every quadratic --- real or complex --- with no guessing?} It exists, and it comes from
a move you already know: \emph{completing the square}, done once on the general equation.
\end{hookbox}

\begin{notesbox}{1. The quadratic formula --- completing the square, done once}
Complete the square on the \emph{general} quadratic $ax^2+bx+c=0$ (same steps as Lesson 2.3, now with
letters). The result is the \textbf{quadratic formula}:
\[ x=\frac{-b\pm\sqrt{\rule{0pt}{1.6ex}\blank{2.4cm}}}{\blank{1.0cm}}. \]
Because completing the square works on \emph{every} quadratic, this formula does too. \textbf{The protocol:}
\begin{enumerate}[leftmargin=1.7em, itemsep=1pt]
  \item Write the equation in \textbf{standard form} $ax^2+bx+c=0$.
  \item Identify $a,b,c$ \emph{with signs}.
  \item Substitute and simplify (radical or $a+bi$).
\end{enumerate}
\textbf{Anchor example.} Solve $x^2-2x-3=0$: \ $a=\blank{0.8cm}$, $b=\blank{0.8cm}$, $c=\blank{0.8cm}$.
\[
\begin{aligned}
  x &= \frac{-(-2)\pm\sqrt{(-2)^2-4(1)(-3)}}{2(1)} \\
    &= \frac{2\pm\sqrt{\blank{1.0cm}}}{2} \\
    &= \frac{2\pm\blank{0.8cm}}{2} \\
    &= \blank{2.2cm}
\end{aligned}
\]
Same zeros as factoring $(x+1)(x-3)$ in 2.2 --- \emph{one method now solves it all.}
\end{notesbox}

\begin{notesbox}{2. The discriminant $b^2-4ac$ --- predict before you solve}
The quantity under the radical, $b^2-4ac$, is the \textbf{discriminant}. Its \emph{sign} alone tells you the
number and type of roots --- \emph{before} you finish solving. Fill in the table and match each to its graph.
\begin{center}
\renewcommand{\arraystretch}{1.4}
\begin{tabularx}{\linewidth}{@{}c l X@{}}
\textbf{If $b^2-4ac$ is\dots} & \textbf{Roots} & \textbf{Graph ($x$-intercepts)} \\
\hline
$>0$ & \blank{4.0cm} & \blank{5.0cm} \\
$=0$ & \blank{4.0cm} & \blank{5.0cm} \\
$<0$ & \blank{4.0cm} & \blank{5.0cm} \\
\end{tabularx}
\end{center}
\vspace{2pt}
\begin{center}
\begin{tikzpicture}[scale=0.36]
  \lgrid{-2}{4}{-5}{3}
  \begin{scope}\clip (-2,-5) rectangle (4,3);
    \draw[very thick, forest, domain=-1.2:3.2, samples=60, smooth] plot (\x,{(\x-1)*(\x-1)-4});
  \end{scope}
  \fill[charcoal] (-1,0) circle (3pt); \fill[charcoal] (3,0) circle (3pt);
  \node[font=\scriptsize, charcoal] at (1,-6.2){$b^2-4ac>0$};
\end{tikzpicture}
\hspace{0.5cm}
\begin{tikzpicture}[scale=0.36]
  \lgrid{0}{4}{-1}{5}
  \begin{scope}\clip (0,-1) rectangle (4,5);
    \draw[very thick, forest, domain=0.05:3.95, samples=60, smooth] plot (\x,{(\x-2)*(\x-2)});
  \end{scope}
  \fill[charcoal] (2,0) circle (3pt);
  \node[font=\scriptsize, charcoal] at (2,-2.4){$b^2-4ac=0$};
\end{tikzpicture}
\hspace{0.5cm}
\begin{tikzpicture}[scale=0.36]
  \lgrid{-3}{3}{-1}{7}
  \begin{scope}\clip (-3,-1) rectangle (3,7);
    \draw[very thick, forest, domain=-2.3:2.3, samples=60, smooth] plot (\x,{\x*\x+1});
  \end{scope}
  \fill[charcoal] (0,1) circle (3pt);
  \node[font=\scriptsize, charcoal] at (0,-2.4){$b^2-4ac<0$};
\end{tikzpicture}
\end{center}
\vspace{-2pt}
Classify \emph{without} solving --- just compute $b^2-4ac$:
\[ x^2-6x+7=0:\ b^2-4ac=\blank{1.6cm}\ \Rightarrow\ \blank{3.2cm}. \]
\[ 4x^2-12x+9=0:\ b^2-4ac=\blank{1.6cm}\ \Rightarrow\ \blank{3.2cm}. \]
\end{notesbox}

\begin{notesbox}{3. When the discriminant is negative --- complex solutions}
A negative discriminant is \emph{not} a dead end. Use $i$ (Lesson 2.4): the roots are a
\textbf{complex-conjugate pair} $a\pm bi$. Solve the 2.4 leftover $x^2+2x+5=0$:
\[ a=1,\ b=2,\ c=5:\quad b^2-4ac=4-20=\blank{1.2cm}. \]
\[
\begin{aligned}
  x &= \frac{-2\pm\sqrt{-16}}{2} \\
    &= \frac{-2\pm\blank{1.0cm}}{2} \\
    &= \blank{2.4cm}
\end{aligned}
\]
These are the \emph{same} solutions completing the square gave in Lesson 2.4 --- now from the formula.
The parabola $y=x^2+2x+5$ never touches the $x$-axis, which is exactly what $b^2-4ac<0$ predicted.
\end{notesbox}

\begin{notesbox}{4. Choose a method \& model}
The formula \emph{always} works, but it is not always fastest. Quick guide:
\begin{center}
\renewcommand{\arraystretch}{1.25}
\begin{tabularx}{\linewidth}{@{}l X@{}}
\textbf{Factoring} & fastest \emph{when} it factors over integers (e.g.\ $x^2-2x-3$). \\
\textbf{Square roots} & for $x^2=k$ or $(x-h)^2=k$ (no middle term). \\
\textbf{Quadratic formula} & the universal fallback --- always works, real or complex. \\
\end{tabularx}
\end{center}
\textbf{Model.} A ball is thrown so its height is $h(t)=-16t^2+24t+6$ (feet, seconds). When does it hit the
ground ($h=0$)? Set up and solve; keep only the physically sensible root.
\begin{work}
  -16t^2+24t+6 &= 0 \\
  t &= \frac{-24\pm\sqrt{24^2-4(-16)(6)}}{2(-16)} \\
    &= \frac{-24\pm\sqrt{960}}{-32} \\
    &= \frac{-24\pm8\sqrt{15}}{-32} \\
    &= \frac{3\pm\sqrt{15}}{4} \\
    &\approx 1.72\text{ s} \quad\text{(reject the negative root)}
\end{work}
\end{notesbox}

\begin{practicebox}
\textbf{Solve with the quadratic formula. Simplify; use $a+bi$ when needed.}
\begin{enumerate}[leftmargin=1.9em, itemsep=6pt]
  \item $x^2+2x-5=0$
    \begin{work}
      x &= \frac{-2\pm\sqrt{2^2-4(1)(-5)}}{2(1)} \\
        &= \frac{-2\pm\sqrt{24}}{2} \\
        &= \frac{-2\pm2\sqrt{6}}{2} \\
        &= -1\pm\sqrt{6}
    \end{work}
  \item $2x^2-4x+3=0$
    \begin{work}
      x &= \frac{4\pm\sqrt{(-4)^2-4(2)(3)}}{2(2)} \\
        &= \frac{4\pm\sqrt{-8}}{4} \\
        &= \frac{4\pm2i\sqrt{2}}{4} \\
        &= 1\pm\frac{\sqrt{2}}{2}i
    \end{work}
  \item Discriminant only: how many/what type of roots does $9x^2-6x+1=0$ have? \blank{3.6cm}
\end{enumerate}
\end{practicebox}

\end{document}
